% !PRINTED 676
\weberpage{676}

%<<CONT
are contained in $S$. For if there is a polygon $\mathfrak{A}_{s+1}$ in $A$
that is not contained in $S$, then

\[
(\mathfrak{A}_1, \mathfrak{A}_2, \ldots, \mathfrak{A}_s, \mathfrak{A}_{s+1})
\]

is a family of dimension $s+1$. Thus it is proved:

7. \gesperrt{The polygons of a class form a family
of a finite dimension $s$. This number shall
be called the dimension of the class.}

If in a class $C$ there are polygons that are divisible by a
polygon $\mathfrak{A}$ of a class $A$, then $C$ is divisible by $A$.
If

\[
C = AB,
\]

then one obtains the dimension of $B$ by selecting from $C$ the
family that is divisible by some polygon $\mathfrak{A}$ in $A$. This
dimension therefore depends only on the two classes $A$ and $C$,
and shall be denoted by $(A, C)$. The symbol for the dimension of
a class $A$ is accordingly $(O, A)$, if $O$ is the class of the
null polygon. We then have

\weq{14}{(A, C) = (A, AB) = (O, B),}

and if $a$ is the order of the class $A$, then by 5.:

\weq{15}{(A, C) \geqq (O, C) - a.}

If we denote by $A$ the proper class which one obtains
when one cancels $\mathfrak{M}$ away everywhere, then the improper class
with the divisor $\mathfrak{M}$ shall be denoted by $\mathfrak{M}A$.

8. \gesperrt{The divisor of an improper class is an
isolated polygon.}

For if $\mathfrak{M}A$ is an improper class, then a polygon $\mathfrak{A}$
can be found in $A$ that is relatively prime to $\mathfrak{M}$. If $\mathfrak{M}'$ is
equivalent to $\mathfrak{M}$, then $\mathfrak{M}'\mathfrak{A}$ is contained in $\mathfrak{M}A$ and consequently
divisible by $\mathfrak{M}$. Hence $\mathfrak{M}'$ too is divisible by $\mathfrak{M}$, and therefore
identical with $\mathfrak{M}$. There is therefore in the class of $\mathfrak{M}$ only the
single polygon $\mathfrak{M}$.

\begin{center}
§ 193. \gesperrt{Normal bases.}
\end{center}

According to the foregoing we can represent a function $z$ of the field
$\Omega$ of the $n$th order in the form

\weq{1}{z = \frac{\mathfrak{U}'}{\mathfrak{U}}}
%CONT>>
